\tema{f) Interpretar los resultados físicos con sus signos}
\paso{1. Tendencia del aluminio.}
Para saber si $F_3$ crece examinamos $F_3'=4.773+4.12x-8.892x^2=4.773+x(4.12-8.892x)$. En el intervalo,
\[
8.892x\leq8.892(0.3872)=3.4429824,\quad
4.12-8.892x\geq0.6770176>0.
\]
Como $x\geq0$, el segundo término no es negativo y $F_3'\geq4.773>0$. Así el logaritmo aumenta. La derivada de $\gamma=e^{F_3}$ es $\gamma'=F_3'e^{F_3}>0$, pues ambos factores son positivos. Por tanto el coeficiente también aumenta.

En cero, $\gamma=e^{-3.8685}\simeq0.02088968$; en el extremo es \GammaFinal. Todos esos valores son menores que 1: la solución presenta desviación negativa respecto de la idealidad en esta escala. La actividad $a=x\gamma$ también aumenta porque $a'=\gamma+x\gamma'>0$.

\paso{2. Tendencia del hierro.}
Gibbs-Duhem da $H_3'=-xF_3'/(1-x)$. Para $x>0$, el numerador $xF_3'$ es positivo y el denominador $1-x$ es positivo; el signo exterior lo hace negativo. Por ello $H_3$ disminuye desde cero y $\gamma_{\Fe}=e^{H_3}$ disminuye desde 1. En el extremo, $H_3=\HierroFinal$ y $\gamma_{\Fe}\simeq0.58867$.

\begin{center}\includegraphics[width=\linewidth]{output/figuras/termodinamica.pdf}\end{center}
\textbf{Figura 3.} A la izquierda, $H_N$; a la derecha, $g_N^E/1000$ en kJ/mol. Cada par usa el mismo orden para Al y Fe. El error máximo de energía frente al cúbico es \ErrorEnergiaLineal\ J/mol para orden 1 y \ErrorEnergiaCuadratico\ J/mol para orden 2. La energía y el coeficiente del Al no tienen por qué clasificar igual la precisión de un truncamiento: combinan las funciones de forma diferente.

\tema{f) Energía negativa, calor y estabilidad: tres conceptos distintos}
\paso{1. Qué dice el signo de $g^E$.}
$g^E=g_{\mathrm{real}}-g_{\mathrm{ideal}}$. Si $g^E<0$, al pasar de la referencia ideal al modelo real disminuye la energía de Gibbs de mezcla. Pero $g^E=h^E-Ts^E$: contiene entalpía y entropía. Por eso $-16.41466$ kJ/mol no se interpreta automáticamente como calor liberado.

\paso{2. Escribir la energía total de mezcla del modelo.}
Sin las contribuciones lineales de los estados puros, que no cambian la segunda derivada,
\[
g_{\mathrm{mezcla}}=RT[x\ln x+(1-x)\ln(1-x)]+g^E.
\]
El primer corchete es la contribución ideal; el segundo término es la corrección de exceso. Para estabilidad local interesa si la curva se dobla hacia arriba, es decir, si $g_{\mathrm{mezcla}}''>0$.

\paso{3. Derivar la parte ideal dos veces.}
Por producto y cadena:
\[
(x\ln x)'=\ln x+1,\qquad
[(1-x)\ln(1-x)]'=-\ln(1-x)-1.
\]
Su suma es $\ln x-\ln(1-x)$. Derivando otra vez: $1/x+1/(1-x)$.

\paso{4. Derivar la parte de exceso con Gibbs-Duhem.}
Definimos $f_N=xP_N+(1-x)H_N$. Por producto:
\[
f_N'=P_N+xP_N'-H_N+(1-x)H_N'.
\]
Gibbs-Duhem dice $xP_N'+(1-x)H_N'=0$; por tanto $f_N'=P_N-H_N$. Derivamos otra vez y sustituimos $H_N'=-xP_N'/(1-x)$:
\[
f_N''=P_N'-H_N'=P_N'+\frac{xP_N'}{1-x}
=P_N'\frac{1-x+x}{1-x}=\frac{P_N'}{1-x}.
\]
\paso{5. Sumar ambas contribuciones para la referencia.}
\[
g_{\mathrm{mezcla}}''=RT\left[\frac1x+\frac1{1-x}+\frac{F_3'}{1-x}\right]>0
\quad\text{para }0<x\leq0.3872.
\]
Todos los factores son positivos en ese intervalo. Esto indica estabilidad local frente a pequeñas fluctuaciones de composición dentro del modelo líquido; no compara globalmente su energía con la de sólidos u otras fases.

\tema{g) Conclusiones, límites y cómo justificar cada afirmación}
\paso{1. Qué permite conocer Taylor.}
La constante describe el límite diluido, la pendiente positiva el aumento inicial de $\ln\gamma_{\Al}$ y las potencias superiores corrigen la forma de esa variación. La expansión permite cuantificar qué se pierde al omitir términos. Esto se sustenta en las derivadas y en la resta explícita de los polinomios.

\paso{2. Calidad obtenida.}
Para el modelo de referencia, los máximos errores relativos en $\gamma_{\Al}$ son 12.7839\,\% con orden 1 y 18.7751\,\% con orden 2. El RMSE del logaritmo sí mejora de orden 1 a 2. La compensación entre el término cuadrático positivo y el cúbico negativo explica esta aparente contradicción. El orden 3 tiene error cero solo porque coincide con la referencia.

\paso{3. Condiciones de uso.}
Con tolerancia de 1\,\% en el coeficiente, el lineal funciona hasta $x\simeq\LimiteLineal$ y el cuadrático hasta $x\simeq\LimiteCuadratico$ en comparación con $F_3$. Para cubrir la malla completa con esa referencia conviene conservar los cuatro términos. Esos umbrales no certifican exactitud experimental.

\paso{4. Dos fuentes diferentes de error.}
Una es simplificar Fe-Al-O como Fe-Al; otra es truncar Taylor dentro del binario. Reducir la segunda no elimina la primera. Además están el error del modelo físico frente a mediciones y los efectos de redondeo numérico. No deben llamarse a todos ``error de Taylor''.

\paso{5. Lo que no se puede inferir.}
Los coeficientes efectivos no identifican de manera única mecanismos atómicos. No se calcula cinética, no se estima $\gamma_{\mathrm O}$ con el binario, no se cambia de temperatura conservando parámetros sin justificación y no se presenta un ajuste experimental independiente. El PDF del artículo no aporta una tabla completa de mediciones EMF con incertidumbres para realizar aquí esa validación.

\paso{6. Por qué no extrapolar a aluminio puro con el redondeo.}
En $x=1$, $F_3(1)=-3.8685+4.773+2.06-2.964=0.0005$. El $H_3$ redondeado contiene $0.001\ln(1-x)$, que diverge al aproximarse a 1. Con $c_2=2.0595$, $\F(1)=0$ y el logaritmo del hierro se cancela. Un cambio diminuto dentro de la malla puede tener consecuencias grandes al extrapolar a una frontera singular. El programa excluye $x\geq1$.

\idea{Se ha desarrollado Taylor de un modelo declarado, se ha medido su truncamiento y se ha mantenido la consistencia entre especies. Ninguna de esas operaciones convierte automáticamente el modelo en una medición de la naturaleza.}
